\subsection{Ordering Guarantees and Misconceptions}

Dictionaries and sets are both hash-based collections, but their visible ordering behavior is not the same. Modern Python dictionaries preserve insertion order. Sets do not provide an insertion-order guarantee.

This creates two common misconceptions:

\begin{itemize}
    \item a dictionary that preserves insertion order is sometimes mistaken for a sorted mapping;
    \item a set that prints in a stable-looking order is sometimes mistaken for an ordered collection.
\end{itemize}

Both interpretations are wrong. A dictionary is still a hash-based key-value mapping. A set is still an unordered uniqueness domain.

\subsubsection{Dictionary Insertion Order: Preserved Order Is Not Sorted Order}

A modern dictionary preserves the order in which keys are first inserted.

\begin{verbatim}
data = {}

data["zebra"] = 1
data["apple"] = 2
data["mango"] = 3

print(f"data: {data!r}")
print(f"type(data): {type(data)}")
print(f"list(data): {list(data)}")
\end{verbatim}

Possible output:

\begin{verbatim}
data: {'zebra': 1, 'apple': 2, 'mango': 3}
type(data): <class 'dict'>
list(data): ['zebra', 'apple', 'mango']
\end{verbatim}

The keys appear as \texttt{"zebra"}, \texttt{"apple"}, and \texttt{"mango"} because that is the insertion order. This is not alphabetical order.

\begin{verbatim}
data = {}

data["zebra"] = 1
data["apple"] = 2
data["mango"] = 3

print(f"list(data):   {list(data)}")
print(f"sorted(data): {sorted(data)}")
\end{verbatim}

Possible output:

\begin{verbatim}
list(data):   ['zebra', 'apple', 'mango']
sorted(data): ['apple', 'mango', 'zebra']
\end{verbatim}

The expression \texttt{list(data)} follows dictionary iteration order. The expression \texttt{sorted(data)} constructs a new sorted list of keys. Sorting is a separate operation.

The dictionary methods expose the same insertion order.

\begin{verbatim}
data = {}

data["zebra"] = 1
data["apple"] = 2
data["mango"] = 3

print(f"keys:   {list(data.keys())}")
print(f"values: {list(data.values())}")
print(f"items:  {list(data.items())}")
\end{verbatim}

Possible output:

\begin{verbatim}
keys:   ['zebra', 'apple', 'mango']
values: [1, 2, 3]
items:  [('zebra', 1), ('apple', 2), ('mango', 3)]
\end{verbatim}

The dictionary is not sorted by key and not sorted by value. It is preserving the order of first insertion.

Updating an existing key changes the value, not the key position.

\begin{verbatim}
quotes = {}

quotes["Homer"] = "D'oh!"
quotes["Jerry"] = "No soup for you!"
quotes["Arthur"] = "It is only a model."

quotes["Jerry"] = "Serenity now!"

print(f"quotes: {quotes!r}")
print(f"list(quotes): {list(quotes)}")
\end{verbatim}

Possible output:

\begin{verbatim}
quotes: {'Homer': "D'oh!", 'Jerry': 'Serenity now!', 'Arthur': 'It is only a model.'}
list(quotes): ['Homer', 'Jerry', 'Arthur']
\end{verbatim}

The key \texttt{"Jerry"} remains in its original position because the key was not newly inserted. Only its associated value changed.

If a key is removed and inserted again, it receives a new insertion position.

\begin{verbatim}
quotes = {
    "Homer": "D'oh!",
    "Jerry": "No soup for you!",
    "Arthur": "It is only a model.",
}

del quotes["Jerry"]
quotes["Jerry"] = "Serenity now!"

print(f"quotes: {quotes!r}")
print(f"list(quotes): {list(quotes)}")
\end{verbatim}

Possible output:

\begin{verbatim}
quotes: {'Homer': "D'oh!", 'Arthur': 'It is only a model.', 'Jerry': 'Serenity now!'}
list(quotes): ['Homer', 'Arthur', 'Jerry']
\end{verbatim}

The second insertion of \texttt{"Jerry"} occurs after the remaining keys, so it appears at the end.

The relevant dictionary methods can be inspected directly.

\begin{verbatim}
data = {"zebra": 1, "apple": 2, "mango": 3}

print(f"type(data): {type(data)}")
print(f"keys in dir:   {'keys' in dir(data)}")
print(f"values in dir: {'values' in dir(data)}")
print(f"items in dir:  {'items' in dir(data)}")
print(f"sort in dir:   {'sort' in dir(data)}")
\end{verbatim}

Possible output:

\begin{verbatim}
type(data): <class 'dict'>
keys in dir:   True
values in dir: True
items in dir:  True
sort in dir:   False
\end{verbatim}

There is no \texttt{sort()} method on a dictionary. Sorting is not part of the dictionary's internal mapping behavior. If sorted output is needed, it must be requested explicitly.

\begin{verbatim}
data = {"zebra": 1, "apple": 2, "mango": 3}

for key in sorted(data):
    print(f"{key!r:<8} -> {data[key]}")
\end{verbatim}

Possible output:

\begin{verbatim}
'apple'  -> 2
'mango'  -> 3
'zebra'  -> 1
\end{verbatim}

This loop uses a sorted list of keys to control traversal. The dictionary itself is not rearranged.

\begin{verbatim}
data = {"zebra": 1, "apple": 2, "mango": 3}

sorted_keys = sorted(data)

print(f"sorted_keys: {sorted_keys}")
print(f"data:        {data}")
\end{verbatim}

Possible output:

\begin{verbatim}
sorted_keys: ['apple', 'mango', 'zebra']
data:        {'zebra': 1, 'apple': 2, 'mango': 3}
\end{verbatim}

The practical rule is:

\begin{quote}
Dictionary order means insertion order, not sorted order.
\end{quote}

\subsubsection{Set Iteration Order: Why Sets Remain Unordered Even When They Appear Stable in Small Examples}

A set is an unordered collection of unique hashable elements. It may print in some order, and that order may appear stable during one run, but the program must not treat that order as meaningful.

\begin{verbatim}
items = {"zebra", "apple", "mango"}

print(f"items: {items!r}")
print(f"type(items): {type(items)}")
print(f"list(items): {list(items)}")
\end{verbatim}

Possible output:

\begin{verbatim}
items: {'apple', 'zebra', 'mango'}
type(items): <class 'set'>
list(items): ['apple', 'zebra', 'mango']
\end{verbatim}

The printed order is not insertion order. It is not sorted order either. It is the current traversal order produced by the set's internal hash-table state.

This is easy to misunderstand because small examples may appear stable.

\begin{verbatim}
items = set()

items.add("zebra")
items.add("apple")
items.add("mango")

print(f"items: {items!r}")
print(f"list(items): {list(items)}")
print(f"sorted(items): {sorted(items)}")
\end{verbatim}

Possible output:

\begin{verbatim}
items: {'apple', 'zebra', 'mango'}
list(items): ['apple', 'zebra', 'mango']
sorted(items): ['apple', 'mango', 'zebra']
\end{verbatim}

The set was filled in the order \texttt{"zebra"}, \texttt{"apple"}, \texttt{"mango"}, but the displayed set does not preserve that insertion order.

The set object exposes membership and mathematical operations, not sequence positioning operations.

\begin{verbatim}
items = {"zebra", "apple", "mango"}

print(f"type(items): {type(items)}")
print(f"__contains__ in dir: {'__contains__' in dir(items)}")
print(f"__iter__ in dir:     {'__iter__' in dir(items)}")
print(f"add in dir:          {'add' in dir(items)}")
print(f"union in dir:        {'union' in dir(items)}")
print(f"sort in dir:         {'sort' in dir(items)}")
print(f"index in dir:        {'index' in dir(items)}")
\end{verbatim}

Possible output:

\begin{verbatim}
type(items): <class 'set'>
__contains__ in dir: True
__iter__ in dir:     True
add in dir:          True
union in dir:        True
sort in dir:         False
index in dir:        False
\end{verbatim}

The absence of \texttt{sort()} and \texttt{index()} is not accidental. A set is not a sequence. It has no promised first element, second element, or last element.

Trying to index a set fails.

\begin{verbatim}
items = {"zebra", "apple", "mango"}

try:
    print(items[0])
except TypeError as err:
    print(f"indexing failed: {err}")
\end{verbatim}

Possible output:

\begin{verbatim}
indexing failed: 'set' object is not subscriptable
\end{verbatim}

If sorted traversal is needed, explicitly sort the set into a list.

\begin{verbatim}
items = {"zebra", "apple", "mango"}

for item in sorted(items):
    print(f"item: {item!r}")
\end{verbatim}

Possible output:

\begin{verbatim}
item: 'apple'
item: 'mango'
item: 'zebra'
\end{verbatim}

The sorted list is ordered. The original set remains unordered.

\begin{verbatim}
items = {"zebra", "apple", "mango"}
ordered_items = sorted(items)

print(f"ordered_items: {ordered_items!r}")
print(f"type(ordered_items): {type(ordered_items)}")
print(f"items: {items!r}")
print(f"type(items): {type(items)}")
\end{verbatim}

Possible output:

\begin{verbatim}
ordered_items: ['apple', 'mango', 'zebra']
type(ordered_items): <class 'list'>
items: {'apple', 'zebra', 'mango'}
type(items): <class 'set'>
\end{verbatim}

The call to \texttt{sorted()} did not transform the set into an ordered set. It created a separate list.

The difference between dictionary and set ordering can be summarized as follows:

\begin{center}
\begin{tabular}{lll}
\textbf{Object} & \textbf{Iteration Behavior} & \textbf{Meaning of Order} \\
\hline
\texttt{dict} & Preserves key insertion order & Order of first key insertion \\
\texttt{set} & No insertion-order guarantee & No semantic order \\
\end{tabular}
\end{center}

A final comparison makes the boundary explicit.

\begin{verbatim}
keys = ["zebra", "apple", "mango"]

data = {}
items = set()

for key in keys:
    data[key] = len(key)
    items.add(key)

print(f"keys source: {keys!r}")
print(f"dict keys:   {list(data)}")
print(f"set values:  {list(items)}")
print(f"sorted set:  {sorted(items)}")
\end{verbatim}

Possible output:

\begin{verbatim}
keys source: ['zebra', 'apple', 'mango']
dict keys:   ['zebra', 'apple', 'mango']
set values:  ['apple', 'zebra', 'mango']
sorted set:  ['apple', 'mango', 'zebra']
\end{verbatim}

The dictionary preserves the order in which keys were inserted. The set stores the same unique values, but it does not preserve their insertion order. Sorting must be requested separately.

The practical rule is:

\begin{quote}
Trust dictionary insertion order. Do not trust set iteration order as program meaning.
\end{quote}